\subsection{\texorpdfstring{Finite \(\Gamma\)-sets and etale algebras}{Finite Gamma-sets and etale algebras}}

Let \(K_s\) be a separable closure of \(K\) (V, p.~45, Prop. 14) and \(\Gamma\) the Galois group of \(K_s\) over \(K\) . By a \(\Gamma\) -set we understand a set \(X\) with an action \((\sigma, x) \mapsto \sigma x\) of the group \(\Gamma\) such that the stabilizer of each point of \(X\) is an open subgroup of \(\Gamma\) . It comes to the same to say that the mapping \((\sigma, x) \mapsto \sigma x\) of \(\Gamma \times X\) into \(X\) is continuous when \(X\) is equipped with the discrete topology.

Let \(\mathbf{X}\) be a finite \(\Gamma\) -set. We define an action of \(\Gamma\) on the K-algebra \(\mathbf{K}_s^{\mathbf{X}}\) of mappings of \(\mathbf{X}\) into \(\mathbf{K}_s\) by the formula

\[
u _ {\sigma} f (x) = \sigma \left(f \left(\sigma^ {- 1} x\right)\right)\tag{12}
\]

for \(\sigma \in \Gamma\) , \(f \in \mathbf{K}_s^X\) and \(x \in \mathbf{X}\) . Let \(\Theta(\mathbf{X})\) be the set of invariants of \(\Gamma\) in \(\mathbf{K}_s^X\) ; this is the sub-K-algebra of \(\mathbf{K}_s^X\) consisting of mappings \(f: \mathbf{X} \to \mathbf{K}_s\) such that \(f(\sigma x) = \sigma(f(x))\) for \(\sigma \in \Gamma\) and \(x \in \mathbf{X}\) .

Lemma 5. --- Let X be a finite \(\Gamma\) -set and let \(x_1, \ldots, x_n\) be points of X such that the orbits \(\Gamma x_1, \ldots, \Gamma x_n\) form a partition of X. For \(1 \leqslant i \leqslant n\) let \(\Delta_i\) be the stabilizer of \(x_i\) in \(\Gamma\) and let \(L_i\) be the field of invariants of \(\Delta_i\) . Then \(L_1, \ldots, L_n\) are separable extensions of finite degree of K, and the mapping \(f \mapsto (f(x_1), \ldots, f(x_n))\) is a K-algebra isomorphism of \(\Theta(X)\) onto \(L_1 \times \cdots \times L_n\) .

By hypothesis the subgroups \(\Delta_{1},\ldots,\Delta_{n}\) of \(\Gamma\) are open and Cor. 5 of V, p.~68, shows that the subextensions \(L_{1},\ldots,L_{n}\) of \(K_{s}\) are of finite degree over K. Clearly they are separable; now the last assertion of Lemma 5 is immediate.

From Lemma 5 and Th. 4 (V, p.~34f.) we obtain immediately the following result.

PROPOSITION 12. --- For every finite \(\Gamma\) -set X the algebra \(\Theta(X)\) is etale over K, of degree equal to the cardinal of X. Moreover, every etale algebra over K is isomorphic to an algebra of the form \(\Theta(X)\) .

Remarks. --- a) It is easy to show that for every K-algebra homomorphism \(\varphi\) of \(\Theta(\mathbf{X})\) into \(\mathbf{K}_s\) there exists a unique element \(x\) of \(\mathbf{X}\) such that \(\varphi(f) = f(x)\) for all \(f \in \Theta(\mathbf{X})\) .

\begin{enumerate}
\def\labelenumi{\arabic{enumi})}
\setcounter{enumi}{1}
\tightlist
\item
  Let X and Y be two finite \(\Gamma\) -sets. Let \(\mathfrak{F}_{\Gamma}(X, Y)\) be the set of mappings \(u\) of X into Y such that \(u(\sigma x) = \sigma u(x)\) for all \(\sigma \in \Gamma\) and all \(x \in X\) . For \(u \in \mathfrak{F}_{\Gamma}(X, Y)\) we define a K-algebra homomorphism \(u^{*}: \Theta(Y) \to \Theta(X)\) by \(u^{*}(f) = f \circ u\) . For every homomorphism \(\Psi\) of \(\Theta(Y)\) into \(\Theta(X)\) there exists a unique element \(u\) of \(\mathfrak{F}_{\Gamma}(X, Y)\) such that \(\Psi = u^{*}\) .
\end{enumerate}

